\section{Momentum smearing and hadronic two-point functions}
\label{sec:wuppertal}
We recapitulate Wuppertal smearing as a generic example of
an iterative smearing algorithm.
Then to clarify notations we discuss the standard construction
of hadronic two-point functions,
before generalizing the smearing by introducing a momentum shift as described
above. The discussion can easily be extended to incorporate generic hadrons and
$n$-point functions with $n\geq 2$. 
We conclude with a suggestion how to
further improve the method. This is relevant for fine lattices
where the iteration count becomes large. In the Appendix
we discuss how an additional Lorentz boost factor can be implemented.
\subsection{Wuppertal smearing}
\label{sec:wup}
The most prominent gauge covariant realization of a smearing function is
Wuppertal smearing~\cite{Gusken:1989ad,Gusken:1989qx}, where
$F=\Phi^n$, with $\Phi$ being defined as
\begin{equation}
(\Phi q)_{\vect{x}} = \frac{1}{1+2d\wupc} \left[q_{\vect{x}}+\wupc \sum_{j=\pm 1}^{\pm d} U_{\vect{x},j} q_{\vect{x}+\unitj}\right]\,. \label{fermsmear}
\end{equation}
Again $d=3$ denotes the dimension of space. The gauge
transporters
\begin{align}\label{eq:gausssmear}
U_{x,-j} = U^{\dagger}_{x-\unitj,j}\,,\quad U_{\vect{x},j}=U_{x,j}\quad\text{with}\,\, x_4\,\,\text{fixed}\,,
\end{align}
can also be spatially APE smeared
links~\cite{Falcioni:1984ei,Alexandrou:1990dq,Bali:2005fu}.
$\unitj$ denotes a vector of length of one lattice unit $a$, pointing
into the direction $j$. In Eq.~\eqref{fermsmear} we suppressed the index $x_4$
since the smearing is diagonal in time. $\wupc$ is a positive
constant and the arbitrary normalization factor $1/(1+2d\wupc)$
is introduced to avoid a numerical overflow for large iteration
counts $n$.

$\Phi$ (and by implication $\Phi^n$) is self-adjoint, a unit matrix
in spinor space and transforms according to the $A_1$ representation of $O_h$
or its little groups.\footnote{$O_h$ symmetry will be reduced to $C_{nv}$
for momenta along a lattice axis ($C_{4v}$), a spatial diagonal ($C_{3v}$)
or a planar diagonal ($C_{2v}$), see, e.g., Ref.~\cite{Hamermesh:100343}.}
The replacement $q\mapsto \Phi^n q$ will therefore not interfere with
the symmetry properties of any interpolator.

The smearing operator $\Phi$ is related to a
discretized covariant Laplacian $\Laplace$:
\begin{align}
\Phi\,q&=q+\frac{\wupc}{1+2d\wupc}a^2\Laplace q\,,\\
a^2\left(\Laplace q\right)_{\vect{x}}
&=-2d\, q_{\vect{x}}+\sum_{j=\pm 1}^{\pm d} U_{\vect{x},j} q_{\vect{x}+\unitj}\,.\label{eq:laplace}
\end{align}
Introducing a fictitious time $\tau=n\Delta\tau$ and defining
$q(\tau)=\Phi^nq(0)$, i.e.\ $q(\tau+\Delta\tau)=\Phi q(\tau)$,
results in the diffusion (or heat) equation
\begin{equation}
\frac{\partial q(\tau)}{\partial\tau}\approx
\frac{q(\tau+\Delta\tau)-q(\tau)}{\Delta\tau}=\alpha\Laplace q(\tau)\,,\label{eq:diffuse}
\end{equation}
where
\begin{equation}
\alpha=\frac{a^2}{\Delta\tau}\frac{\wupc}{1+2d\wupc}\,.
\label{eq:alphadef}
\end{equation}
This is solved by $q(\tau)\approx \exp\left(\alpha\tau\Laplace\right)q(0)$.
Starting from $\delta$-sources in position and colour space
$q_{\vect{x}}^a(0)=\delta_{\vect{x}\vect{0}}\delta^{ab}$,
assuming the free case $U_{\vect{x},j}=\mathds{1}_3$
and large distances $r=|\vect{x}|\gg a$, Eq.~\eqref{eq:diffuse}
is obviously solved by a Gaussian with
\begin{equation}
\sigma(\tau)=\sqrt{2\alpha\tau}=\sqrt{2na^2}\sqrt{\frac{\wupc}{1+2d\wupc}}\label{eq:width}
\end{equation}
being the square root of the variance,
i.e.\ the smearing corresponds to the smearing kernel Eq.~\eqref{eq:smgauss}.
Employing a parallel transporter within the covariant Laplacian
Eqs.~\eqref{eq:gausssmear} and \eqref{eq:laplace}
that is close to unity, like
$d$ dimensional APE smeared gauge links,
means that this
Gaussian shape is a good approximation, see, e.g., Ref.~\cite{Bali:2011rd}.

The diffusivity $\alpha$ obviously is maximal for large values of the
parameter $\wupc$
($\alpha\rightarrow \frac16 a^2/\Delta\tau$ for $\wupc\rightarrow\infty$).
For small values ($\alpha\approx \wupc\, a^2/\Delta\tau$),
the iteration count
$n$ to achieve a given smearing radius\footnote{Note that the root mean
squared radius in $d$ dimensions reads
$\sqrt{d}\sigma$. Also
note that the width of the gauge invariant combination
$q^{\dagger}_{\vect{x}}q_{\vect{x}}$ Eq.~\eqref{density}
is smaller by a factor $\sqrt{2}$.}
$\sqrt{3}\sigma$ is larger but the
resulting wave function will be smoother. Equation~\eqref{eq:width}
highlights that, to keep the smearing radius fixed in physical
units, the iteration
count needs to be increased in proportion to the square of the
inverse lattice spacing. Obviously,
this becomes computer time intensive towards the
continuum limit. We remark that at small lattice spacings
the diffusion equation~\eqref{eq:diffuse} should be
solved using a smarter method than iteratively applying the smearing operator
$\Phi$ defined in Eq.~\eqref{eq:gausssmear}. We will return to this
in Sec.~\ref{sec:improve}.

\subsection{Construction of two-point functions} 
As an example we discuss the construction of momentum projected
pion and nucleon two-point functions:
\begin{align}
C_{\pi}(\vect{p},t)
&=\langle \pi_{\vect{p}} (t) \pi^{\dagger}_{\vect{p}} (0)\rangle\,,\nonumber\\
C_{N}(\vect{p},t)&=
\langle N_{\vect{p}}^{\alpha} (t) \overline{N}_{\vect{p}}^{\beta}(0)\rangle \mathbb{P}^{\alpha\beta}\,,
 \label{interpol}
\end{align}
where $\mathbb{P}=\frac{1}{2}(\mathds{1}+\gamma_4)$ denotes a projector onto
positive parity.\footnote{Note that for $\vect{p}\neq\vect{0}$
this projector will not completely remove the negative parity contribution.
To improve on this, instead one can employ the oblique
projector~\cite{Lee:1998cx}
$\mathbb{P}_{\vect{p}}=\frac{1}{2}\{\mathds{1}+[m_{N^*}/E_{N^*}(\vect{p})]\gamma_4\}$,
where $m_{N^*}$ and $E_N^*$ denote mass and energy of the
nucleon's negative parity partner. Since this state is higher
in mass than the nucleon, for simplicity, here we abstain from attempting this.}
Without smearing the pion and nucleon interpolators
are local quark bilinears and trilinears:
\begin{align}
\pi_{\vect{p}}= a^3\sum_{\vect{x}}e^{i\vect{p}\cdot\vect{x}} \bar{u}_{\vect{x}}\gamma_5 d_{\vect{x}}\,,\\
N^{\alpha}_{\vect{p}} = a^3\sum_{\vect{x}} e^{i\vect{p}\cdot\vect{x}} \left[ u^t_{\vect{x}} C \gamma_5 d_{\vect{x}} \right]u^{\alpha}_{\vect{x}}\,,
\end{align}
where $C$ is the charge conjugation operator and $u_{\vect{x}}$ and
$d_{\vect{x}}$ annihilate up and down quarks, respectively, at the
spatial position $\vect{x}$.

For the pion the Wick contraction then gives
\begin{align}
C_{\pi}(\vect{p},t)
&= -L^3a^3\sum_{\vect{x}}
e^{i \vect{p}\cdot\vect{x}} \left\langle\bar{u}_x \gamma_5 d_{x} d^{\dagger}_0 \gamma_5
\bar{u}^{\dagger}_0\right\rangle\label{eq:contract} \\\nonumber
&=-L^3a^3\sum_{\vect{x}}\left\langle\tr\left[(d\bar{d})_{x0}\gamma_5
(u\bar{u})_{0x}\gamma_5\right]\right\rangle e^{i \vect{p}\cdot\vect{x}}\\
&=-L^3a^3\left\langle\sum_{\vect{x}}\tr\left(M^{-1\dagger}_{x0}M^{-1}_{x0}
\right)e^{i \vect{p}\cdot\vect{x}}\right\rangle\,,\nonumber
\end{align}
where $x=(\vect{x},t)$ and the trace is over spin and colour.
Momentum projection at the source is not
necessary, due to the translational symmetry of expectation values
and $L^3=N^3a^3$ denotes the three-volume that corresponds to this omitted sum.
$M$ is a Wilson-like lattice Dirac matrix
with the quark mass $m_u=m_d$ and we
have used\footnote{Note that the assignment
$\bar{d}=d^{\dagger}\gamma_4$ is convenient but
not consistent in Euclidean spacetime, where
the Dirac operator is $\gamma_5$- instead of
$(\gamma_4=\gamma_0)$-Hermitian.
This is the reason for the negative sign in Eq.~\eqref{eq:contract}.
In the simulation, we correct for this phase and obtain a
positive two-point function.} $d^\dagger\gamma_5\bar{u}^\dagger=d^\dagger
\gamma_5\gamma_4u=-\bar{d}\gamma_5u$, the Grassmann nature of the
quark field creation and annihilation operators and replaced
$\langle d\bar{d}\rangle_d=\langle u\bar{u}\rangle_u=M^{-1}$,
where the subscripts $d$, $u$ denote integrating out the respective
Grassmann field
on a given gauge background. In the last step we also
exploited $\gamma_5$-Hermiticity:
$\gamma_5M^{-1}\gamma_5={M^{-1}}^{\dagger}$.
For the nucleon one can easily work out an analogous expression that contains
$\epsilon^{abc}\epsilon^{a'b'c'}(M^{-1})_{x0}^{aa'}
(M^{-1})_{x0}^{bb'}(M^{-1})_{x0}^{cc'}$, where we have suppressed the
spinor indices and $a,b,c,a',b',c'\in\{1,2,3\}$ run over fundamental
colour.
Equation~\eqref{eq:contract} can now be evaluated, solving the linear systems
\begin{equation}
\sum_{x,\alpha, a}M^{\beta' b'\alpha a}_{x'x}S^{\alpha a \beta  b}_x=\delta_{x'0}\delta^{b\beta'\beta}\delta^{b'b}\label{eq:prop}
\end{equation}
for twelve $\delta$-sources
(for each source spin $\beta$ and colour $b$) to obtain the point-to-all
propagator $S$ with sink position, spin and colour indices
$x$, $\alpha$ and $a$, respectively. Then,
\begin{equation}
C_{\pi}(\vect{p},t)\propto\left\langle\sum_{\vect{x}}\tr\left(S^{\dagger}_xS_x\right)e^{i\vect{p}\cdot\vect{x}}
\right\rangle\,,\label{eq:picontract}
\end{equation}
where again $x_4=t$. The construction of the analogous nucleon
two-point function from three point-to-all propagators
is straightforward. Note that in that case no Hermitian adjoint
will appear.

Smearing is diagonal in Euclidean time (hence
we suppressed the $t$ dependence) and trivial in Dirac spin. 
So, obviously, the smearing operator $F=\Phi^n$ commutes with
any $\Gamma$ structure, however,
it does not commute with covariant derivatives that may appear
within the hadronic interpolator. 
Smearing can easily be implemented
at the source, replacing $S=M^{-1}\delta$
by $S^F=M^{-1}F\delta$, see Eq.~\eqref{eq:prop}.
Note that, due to the fact that
$F$ is a unit matrix in spinor space, 
the same smeared source can be employed for all four
spinor components.
To obtain a so-called
smeared-smeared two-point function, the argument of the sum
in Eq.~\eqref{eq:picontract} will usually be replaced by
$\tr[(FS^{F})^{\dagger}_{x}(FS^{F})_{x}]e^{i\vect{p}\cdot\vect{x}}$:
every source smearing requires new inversions while
sink smearing needs to be carried out on all time slices
of interest.\\

We remark that momentum sources have been used for quite some time, see, e.g.,
Refs.~\cite{Gockeler:1998ye,Brown:2012tm,Bali:2015msa}.
Injecting momentum into quark sources is
necessary (and has been done) in the context
of the one-end-trick of Refs.~\cite{Sommer:1994gg,Foster:1998wu},
where one usually employs colour diagonal (complex) $\mathbb{Z}_N$ or
$\textmd{U}(1)$ random sources. Generalizing this
to $\textmd{SU}(3)$ noise is in fact equivalent to using a non-gauge fixed
wall source~\cite{Billoire:1985yn}, which does not change expectation
values. Momentum was however also injected into gauge fixed wall sources
in Refs.~\cite{Detmold:2012wc,Basak:2012py}
(``color wave propagator''), favourably affecting not only
statistical errors but also ground state overlaps. While these latter
references share some of our motivation, the method presented
here is quite different. For instance, in our case
the hadron's total momentum still needs
to be injected at the source and $\vect{k}$ is not quantized.\\

Finally, we remark that the asymmetry of only carrying
out the position sum at the sink often is exploited
to reduce the statistical errors of heavy-light meson
correlation functions, by only smearing the heavy quark with
$F^2=FF$ at the source, instead of smearing each quark with $F$.
Clearly, it would be advantageous if for each momentum
smearing parameter of interest only the heavy quark propagator had
to be recomputed. Unfortunately, momentum smearing (as well as
traditional smearing) as
an operation will not commute with momentum projection, unless
$\vect{p}=\vect{0}$. Therefore, a different distribution of the smearing
between the quarks will result in a different
and not necessarily optimal ground state overlap.

\subsection{Momentum (Wuppertal) smearing}
\label{sec:newsmear}
As explained in Sec.~\ref{sec:basic},
if we inject a momentum $\vect{p}$ into the hadron, it may be
a good idea to distribute at least part of it among the quarks.
Their smearing functions should therefore ideally be centred around
a value $\vect{k}\neq\vect{0}$, where $\vect{k}$ is
some fraction of $\vect{p}$, to best resemble the
wave function of the physical state we wish to create.\\

The Fourier transform of a Gaussian,
centred about $\vect{k}$  in momentum space, reads
\begin{equation}
f^{\sigma}_{(\vect{k})\,\vect{x}-\vect{y}}\propto
\exp\left[-\frac{(\vect{x}-\vect{y})^2}{2\sigma^2}-i\vect{k}\cdot(\vect{x}-\vect{y})\right]\,.\label{eq:momsf}
\end{equation}
Similarly, momentum can be injected also into differently
shaped smearing functions, a possibility that we shall not explore
here.  $q(\tau)=F_{(\vect{k})}^{\sigma(\tau)}q(0)$
solves the heat equation with a constant drift term
\begin{equation}
\label{eq:drift}
\frac{\partial q(\tau)}{\partial\tau}=
\alpha\left(\vect{\nabla}+i\vect{k}\right)^2 q(\tau)\,.
\end{equation}

In analogy to the discussion of Sec.~\ref{sec:wup},
in the free case the above smearing function can iteratively be constructed
from ``momentum'' Gau\ss{} smearing steps,
$F_{(\vect{k})} =\Phi_{(\vect{k})}^n$, introducing a phase 
into Eq.~\eqref{fermsmear}:
\begin{equation}
(\Phi_{(\vect{k})} q)_{\vect{x}} = \frac{1}{1+2d\wupc} \left[q_{\vect{x}}+\wupc \sum_{j=\pm 1}^{\pm d} U_{\vect{x},j}e^{i\vect{k}\cdot\unitj} q_{\vect{x}+\unitj}\right]\,, \label{fermsmear2}
\end{equation}
where $\Phi_{(\vect{0})}=\Phi$.
One can easily show that the variance $\sigma^2(\tau)$ is still given
as in Eq.~\eqref{eq:width}. Moreover, as expected
$\Phi_{(\vect{k})}$ remains self-adjoint:


\begin{align}
\Phi_{(\vect{k})\vect{x}\vect{y}}&=\frac{1}{1+2d\wupc}\left\{\delta_{\vect{x},\vect{y}}+
\wupc \sum_{j=1}^{3} \left[\delta_{\vect{x}+\unitj,\vect{y}}
U_{\vect{x},j} e^{i\vect{k}\cdot\unitj}
+\delta_{\vect{y}+\unitj,\vect{x}}U^{\dagger}_{\vect{y},j} e^{-i\vect{k}\cdot\unitj}
\right]\right\}\,,\nonumber\\
\Phi_{(\vect{k})\vect{y}\vect{x}}^{\dagger}
&=
\frac{1}{1+2d\wupc}\left\{\delta_{\vect{y},\vect{x}}+
\wupc \sum_{j= 1}^{3} \left[\delta_{\vect{y}+\unitj,\vect{x}}
U_{\vect{y},j}^{\dagger} e^{-i\vect{k}\cdot\unitj}
+\delta_{\vect{x}+\unitj,\vect{y}}U_{\vect{x},j} e^{i\vect{k}\cdot\unitj}
\right]\right\}
=\Phi_{(\vect{k})\vect{x}\vect{y}}\,.\label{eq:transpose}
\end{align}


We will use the same smearing for quarks and
antiquarks. For the $\pi^-$ meson this amounts to the replacements
$d\mapsto \Phi^n_{(\vect{k})}d$,
$u\mapsto \Phi^n_{(-\vect{k})}u$, $\bar{d}\mapsto \bar{d}\Phi^n_{(\vect{k})}$
and $\bar{u}\mapsto \bar{u}\Phi^n_{(-\vect{k})}$
within Eq.~\eqref{eq:contract}. Note that the $\Phi^n_{(-\vect{k})}$ smearing is needed
as in the contraction
with the momentum projector, $(\delta_{\vect{x}\vect{x}'}e^{i\vect{p}\cdot\vect{x}})
\Phi^n_{(\vect{k})\vect{x}\vect{y}}\Phi^n_{(\vect{k})\vect{x}'\vect{y}'}$,
transposing the ordering of the indices of the first smearing operator
gives $\Phi^n_{(\vect{k})\vect{x}\vect{y}}=\Phi^n_{(-\vect{k})\vect{y}\vect{x}}$
in the free case,
see Eq.~\eqref{eq:freec}.
Exploiting the Hermiticity of $\Phi^n_{\vect{k}}$, we then obtain
\begin{align}
C_{\pi}(t)
\propto\sum_{\vect{x}}&\left\langle\!\tr\!\left\{\left[\left(\Phi^n_{(-\vect{k})}M^{-1}\Phi^n_{(-\vect{k})}\right)_{x0}\right]^{\dagger}\right.\right.\nonumber\\
&\qquad\times\left.\left.\left(\Phi^n_{(\vect{k})}M^{-1}\Phi^n_{(\vect{k})}\right)_{x0}
\right\}\right\rangle e^{i \vect{p}\cdot\vect{x}}\,.
\end{align}
In contrast, for baryonic two-point functions, where
all quarks propagate in the forward direction, only $\Phi^n_{(+\vect{k})}$
will appear. 
The pion two-point function can now be constructed in analogy
to Eq.~\eqref{eq:picontract} from
\begin{equation}
C_{\pi}(\vect{p},t)\propto\left\langle\!\sum_{\vect{x}}\tr\!\left[
\left(\Phi^n_{(-\vect{k})}S_{(-\vect{k})}\right)^{\!\dagger}_{\!x}\!\!\left(
\Phi^n_{(\vect{k})}S_{(\vect{k})}\right)_{\!x}\right]\!e^{i\vect{p}\cdot\vect{x}}\!
\right\rangle\,,
\end{equation}
where the source-smeared
point-to-all propagator is defined as [see Eq.~\eqref{eq:prop}]
\begin{equation}
{S_{(\vect{k})}}^{\alpha a \beta  b}_x=
\sum_{x',\alpha',a',z,c}\!\!\!\!\!\left(M^{-1}\right)^{\alpha a\alpha' a'}_{x,x'}\!\!\delta^{\alpha'\beta}
{\Phi^n_{(\vect{k})}}_{x'z}^{a'c}
\delta_{z0}\delta^{cb}\,,
\end{equation}
or, in shorthand notation: $S_{(\vect{k})}=M^{-1}\Phi_{(\vect{k})}^n\delta$.
For mesons the above twelve linear systems need to be solved, both for
$\delta$-sources smeared with $\Phi_{(\vect{k})}^n$ and
with $\Phi_{(-\vect{k})}^n$, while for baryons smearing with
$\Phi_{(\vect{k})}^n$ is sufficient.

We remark that the momentum smearing Eq.~\eqref{fermsmear2} can very
easily be implemented by substituting the (APE smeared) transporters
$U_{\vect{x},j}\mapsto U_{\vect{x},j}e^{i\vect{k}\cdot\unitj}$
and then iterating the usual Wuppertal smearing
Eq.~\eqref{fermsmear} on these modified
$\textmd{U}(3)$ links.

\subsection{Free field investigation}
\label{freefield}
Having defined the smearing and how the contractions are to be
carried out, we are now in the position to address the question
what value of $\vect{k}$
should be chosen. Naively, one may expect $\vect{k}\approx \vect{p}/2$
and $\vect{k}\approx\vect{p}/3$ to be optimal for mesons and baryons,
respectively, that are composed of degenerate quarks and carry
a total momentum
$\vect{p}$. This is indeed what we will find here for
the non-interacting case.
The interacting case somewhat deviates from this as we will
see in Sec.~\ref{sec:results} below.

Setting $U_{\vect{x},j}=\mathds{1}$ within Eq.~\eqref{fermsmear2}
gives the free case smearing function Eq.~\eqref{eq:momsf}
for a large volume $L^3\gg\sigma^3$ and number of smearing iterations
so that $\sigma\gg a$.
Smearing the quark and antiquark annihilation operators
at the sink with $F_{(\vect{k})}$
and $F_{(-\vect{k})}$, respectively, we obtain
the momentum projected smeared pion interpolator,
\begin{align}
\pi_{\vect{p}}&\propto
\sum_{\vect{x}, \vect{y},\vect{y}^{\prime}}
\bar{u}_{\vect{y}'}f_{(-\vect{k})\vect{y}'-\vect{x}} e^{i\vect{p}\cdot\vect{x}}
f_{(\vect{k})\vect{x}-\vect{y}}\gamma_5d_{\vect{y}}\nonumber\\
&\propto\sum_{\vect{x}, \vect{y},\vect{y}^{\prime}}\exp\left[-\frac{(\vect{y}-\vect{x})^2+(\vect{y}'-\vect{x})^2}{2\sigma^2}\right]
 e^{i\vect{p}\cdot\vect{x}}e^{-2i\vect{k}\cdot\vect{x}} e^{i\vect{k}\cdot(\vect{y}+\vect{y}^{\prime})} \left[\bar{u}_{\vect{y}^{\prime}} \gamma_5 d_{\vect{y}}\right]\label{eq:freec}\\\nonumber
&\propto \exp\left[-\frac{\sigma^2}{2}(\vect{p}-2 \vect{k})^2\right]
\sum_{\vect{Z},\vect{\Delta}}
\exp\left[-\frac{\vect{\Delta}^2}{\sigma^2}\right]
e^{i\vect{p}\cdot \vect{Z}} \left[\bar{u}_{\vect{Z}+\vect{\Delta}} \gamma_5 d_{\vect{Z}-\vect{\Delta}}\right]\,,
\end{align}

where we have defined centre and relative coordinates
\begin{align}
\vect{Z} = \frac{\vect{y}+\vect{y}^{\prime}}{2}\,,\quad
\vect{\Delta} =\frac{\vect{y}^{\prime}-\vect{y}}{2}\,.
\end{align}
Note that the components of $\vect{Z}/a$ and $\vect{\Delta}/a$
can be integer or half-integer valued, subject to
the constraints $(Z_j+\Delta_j)/a\in\mathbb{Z}$.
The result of Eq.~\eqref{eq:freec} is indeed maximized for
$\vect{k}=\vect{p}/2$:
in the free case $\vect{k}=\vect{p}/2$ is
the optimal smearing parameter for mesons
with mass-degenerate valence quarks while we encounter an exponential
suppression in $\vect{p}^2$ for the conventional
$\vect{k}=\vect{0}$ smearing. The broader the wave function in coordinate
space, the more
important becomes the correct choice of $\vect{k}$ as one can
see from the first exponent in Eq.~\eqref{eq:freec} (as well as from
Eq.~\eqref{kerfour}).
Unsurprisingly, for baryonic interpolators $\vect{k}=\vect{p}/3$
would be the optimal choice.

We have demonstrated that in the free case $\vect{k}$ can be
interpreted as the momentum
carried by the smeared quark. If the quarks differ in mass
or, like for the example of the nucleon, the interpolator
is not symmetric with respect to the quark flavour,
injecting different $\vect{k}$-values into different quark fields
may be advisable. In the interacting case
the interpretation is not as straightforward:
neither is the interpolator directly related to any of the
usual definitions of a wave function nor will all
momentum be carried by the valence quarks.

\subsection{Non-iterative (momentum) smearing}
\label{sec:improve}
We propose replacing iterative smearing, where the iteration count
diverges with $a^{-2}$ towards the continuum limit,
by a more refined method.
In this context we show how
to introduce phase factors (and shape functions) into
the ``distillation'' (or Laplacian-Heaviside) method of
Ref.~\cite{Peardon:2009gh}.
Although we already present the basics of the method here, systematic
tests are yet to be completed.
Another natural extension,
which we will investigate numerically in Sec.~\ref{sec:results}, is to
Lorentz boost the smearing function.

We define eigenvectors $|v_{\ell}\rangle$ and eigenvalues $-\omega_{\ell}^2$
of a covariant Laplacian with (smeared) gauge transporters,
at a fixed Euclidean time:
\begin{equation}
\label{eq:eigen}
\left(\Laplace +\omega_\ell^2\right)|v_{\ell}\rangle=0\,.
\end{equation}
Since the Laplacian is self-adjoint, the bra-ket notation is convenient
in the present context.
For a lattice of $N^d$ points per time slice ($L=Na$), $3N^d$ 
linearly independent eigenvectors $|v_{\ell}\rangle$ with
components $v^{a}_{\ell\vect{x}}$ exist.
Such eigenvectors were for instance used in Ref.~\cite{Peardon:2009gh}.
The heat equation~\eqref{eq:diffuse} is solved by
\begin{align}
|q(\tau)\rangle&=\exp\left(\alpha\tau\Laplace\right)|q(0)\rangle\nonumber\\
&=\sum_{\ell}| v_{\ell}\rangle\exp\left(-\alpha\tau\omega_{\ell}^2\right)\langle v_{\ell}|q(0)\rangle
\,.\label{eq:noniter}
\end{align}
If we start from a $\delta$-function at $\tau=0$, this results in a Gaussian of
variance $\sigma^2=2\alpha\tau$.

To implement momentum smearing one
can easily replace $U_{\vect{x},j}\mapsto U_{\vect{x},j}e^{i\vect{k}\cdot\unitj}$
within the covariant Laplacian $\Laplace\mapsto\Laplace_{(\vect{k})}\sim
(\vect{\nabla}+i\vect{k})^2$ and then
recompute the eigenvectors and -values:
\begin{equation}
\label{eq:eigen2}
\left(\Laplace_{(\vect{k})} +\omega_{(\vect{k})\ell}^2\right)|v_{(\vect{k})\ell}\rangle=0\,.
\end{equation}
The natural generalization of Eq.~\eqref{eq:noniter} for moving
particles then reads:
\begin{equation}
|q_{(\vect{k})}(\tau)\rangle
=\sum_{\ell}| v_{(\vect{k})\ell}\rangle\exp\left(-\alpha\tau\omega_{(\vect{k})\ell}^2\right)\langle v_{(\vect{k})\ell}|q(0)\rangle
\,.
\label{eq:noniter2}
\end{equation}

The motivation for computing eigenvectors of the covariant
Laplacian is to truncate Eqs.~\eqref{eq:noniter} or \eqref{eq:noniter2}
at a value $\ell_{\max}$ where $\omega_{\ell_{\max}+1}^2>\omega_{\max}^2$.
As $\alpha\tau=\sigma^2/2$, achieving a suppression by a
factor $e^{-2}$ requires $\omega_{\max}\approx 2/\sigma$.
In general, the wider the smearing function in coordinate space
the less eigenvectors will be needed. The extreme opposite
limit $\alpha\tau=0\Rightarrow\sigma=0$ corresponds to the Laplacian-Heaviside
method proposed in Ref.~\cite{Peardon:2009gh}, where summing over all
eigenvectors will ultimately result in
a $\delta$-function in position space while truncating 
at some finite $\ell_{\max}$ value
gives a bell shape~\cite{Peardon:2009gh}. (In the free case
the modulus would be a sum of sines.) The same holds for
the sources suggested in Ref.~\cite{Brown:2012tm} that correspond to
sums of eigenvectors of the non-interacting Laplacian.

It is trivial to work out more details in the free case, where
one basically encounters a one dimensional problem.
For instance, setting $\omega= j (2\pi/L)\gtrsim 2/\sigma$,
where $j\in\mathbb{Z}\setminus\{0\}$
means
$|j|\gtrsim L/(\pi\sigma)$. Therefore, in $d$ dimensions
$\ell_{\max}\gtrsim 3\times[\pi^{d/2}/\Gamma(d/2+1)]\times[L/(\pi\sigma)]^d$, where the
factor $3$ is due to colour and the next factor is the volume
of a $d$ dimensional unit-sphere ($4\pi/3$ for $d=3$): the number of
required eigenvectors increases in proportion to the
spatial volume in physical units
but is independent of the lattice spacing $a$.
In contrast, in the case of iterative smearing, reducing the lattice spacing
at a fixed value of $\sigma$ increases the iteration count
in proportion to $a^{-2}$, independent of the volume.
For our smearing size $\sigma\approx0.45\units{fm}$ and a spatial volume
$L^3=(6\units{fm})^3$, which
would ensure $Lm_{\pi}>4$ even at the physical mass point, we obtain
$\ell_{\max}\approx 960$ while
for a $(3\,\units{fm})^3$ box
about 120 eigenvectors should suffice.
Therefore, satisfactory results in terms of
ground state overlaps appear to be within reach, employing moderately
large numbers of eigenvectors. This is at present under investigation.

Non-Gaussian shapes can easily
be modelled too, e.g., by multiplying in a ``free form'' weight
function, $\left(|v_{\ell}\rangle\langle v_{\ell}|\right)_{\vect{x}\vect{y}}
\mapsto \sum_{\vect{xy}}h_{\vect{x}\vect{y}}\exp\left[\alpha\tau(\vect{x}-\vect{y})^2\right]\left(|v_{\ell}\rangle\langle v_{\ell}|\right)_{\vect{x}\vect{y}}$,
in analogy to Ref.~\cite{vonHippel:2013yfa}.
However, the numerical complexity of this operation is
$\mathcal{O}(N^{2d})$ for each time slice and eigenvector. At the
source this can be reduced to $\mathcal{O}(N^{d})$
if the solution is only required
for a fixed $\delta$-source position
$|q(0)\rangle=|\delta_{\vect{y}_0}^a\rangle$
($\langle\vect{y},b|\delta_{\vect{y}_0}^a\rangle=(|\delta_{\vect{y}_0}^a\rangle)_{\vect{y}}^b=\delta_{\vect{y}\vect{y}_0}\delta^{ba}$).
Other, numerically less expensive options include
placing non-linear functions of $\omega^2_{\ell}$ in the exponents
of Eqs.~\eqref{eq:noniter} and \eqref{eq:noniter2}, employing
linear combinations of Gaussians or substituting the covariant
Laplacian by a different
operator within Eqs.~\eqref{eq:eigen} and \eqref{eq:eigen2}, e.g.,
introducing an anisotropy.

Once the eigenvectors Eq.~\eqref{eq:eigen} or
Eq.~\eqref{eq:eigen2} of the APE or HYP
smeared covariant Laplacian have been computed,
any smearing radius can be realized, replacing
$\alpha\tau$ by the targeted $\sigma^2/2$ value within the exponent of
Eq.~\eqref{eq:noniter} or Eq.~\eqref{eq:noniter2}.
Another potential advantage of the non-iterative smearing is
that, instead of solving for smeared sources
$|q_{(\vect{k})}(\tau)\rangle$ that have evolved from
a $\delta$-source $|q(0)\rangle=|\delta_{\vect{y}_0a}\rangle$,
one can also directly apply the inverse lattice Dirac operator to
the eigenvectors, thereby
constructing so-called perambulators~\cite{Peardon:2009gh}:
\begin{equation}
{\tau_{(\vect{k})}}^{\alpha\beta}_{m t\, \ell 0}=\left\langle v_{(\vect{k})m t}\left|\left(M^{-1}\right)^{\alpha\beta}_{t0}\right|v_{(\vect{k})\ell 0}\right\rangle\,.
\end{equation}
The inner product above is over spatial position and colour,
replacing these indices by the eigenvector labels
$\ell$ at the sink and $m$ at the source. The respective
Euclidean times are
explicitly shown as additional eigenvector (and perambulator) subscripts.
These perambulators can then be folded with the appropriate
eigenvalue Gaussians during the construction
of hadronic $n$-point functions.

We abstain from repeating here the steps outlined in
Ref.~\cite{Peardon:2009gh} how to construct hadronic two-point
functions from these perambulators as there is only one
difference: In the contractions over $\ell$ and $m$ weight factors
$\exp(-\sigma^2\omega_{(\vect{k})\ell 0}^2/2)$
and $\exp(-\sigma^2\omega_{(\vect{k})m t}^2/2)$ should be introduced,
where $\omega_{(\vect{k})m t}^2$ refers to the $m$th eigenvalue
of the negative covariant Laplacian defined on time slice $t$.

On one hand using perambulators will require a larger number of inversions
and computationally more expensive contractions
than starting from $\delta$-sources.
On the other hand, due to volume averaging, the statistical
errors of the perambulator method will be smaller and this method allows
for more flexibility in the subsequent construction of hadronic
$n$-point functions.
Whether the use of perambulators
or of traditional (smeared) point-to-all propagators is preferable
will therefore depend on the problem at hand and in particular on
the number of different hadrons we are interested in. We have
shown that the same quark smearing can be employed in both cases.

Finally, we remark that in certain situations recomputing the eigenvectors
for different values of $\vect{k}$
may be avoidable. In the interacting case, gauge covariant derivatives
$\nabla^2_j$ will depend on all spatial coordinates, including
$x_i$ with $i\neq j$, and will therefore not commute with each other.
However, our intuition was based on the free case and
also the spatially APE~\cite{Falcioni:1984ei}
or HYP~\cite{Hasenfratz:2001hp} smeared gauge covariant
transporters are close to unity.
Assuming translational invariance, the coordinates can be separated and
the components of an eigenvector read
$v^a_{\vect{x}}\propto\sin(\omega_1^ax_1)\cdots\sin(\omega_d^ax_d)$,
where $\omega_j^a=2\pi m_j^a/L$ and $m_j^a$ are integer valued.
The corresponding eigenvalue of the negative Laplacian
is given as $\omega^2=\vect{\omega}^2$ where
the frequencies for different colour components
are constrained: $\vect{\omega}^2:={\vect{\omega}^1}^2={\vect{\omega}^2}^2={\vect{\omega}^3}^2$. Defining
\begin{equation}
|\tilde{v}_{(\vect{k})\ell}\rangle=e^{-i\vect{k}\cdot\vect{x}}|v_{\ell}\rangle
\label{eq:phases}
\end{equation}
then gives
$\left[\left(\vect{\nabla}+i\vect{k}\right)^2+\omega_{\ell}^2\right]|\tilde{v}_{(\vect{k})\ell}\rangle=0$.

In the interacting case this will not exactly solve the heat equation
with drift~\eqref{eq:drift} but
the approximation $\omega_{(\vect{k})\ell}\approx\omega_{\ell}$,
$|v_{(\vect{k})\ell}\rangle\approx|\tilde{v}_{(\vect{k})\ell}\rangle$ should
be sufficient to construct a gauge covariant Gaussian shape
with the intended phase factors.
Note that the phases Eq.~\eqref{eq:phases} appear both
in the bra- and in the ket-vector of Eq.~\eqref{eq:noniter2},
such that only relative
phases between two spatial positions matter and the choice of the zero
point becomes irrelevant. The clear disadvantage of the explicit
multiplication by phase factors Eq.~\eqref{eq:phases} is that these
have to obey the lattice periodicity, i.e.\ $k_j\in (2\pi/L)\mathbb{Z}$.
However, such a restriction may be tolerable on large lattices.
Also introducing twisted fermionic boundary
conditions~\cite{Bedaque:2004kc,Sachrajda:2004mi} may
provide a way to increase the flexibility of the choice of $\vect{k}$.

